\section{Conclusion}

Sparse Readout Prism factorizes the final readout into a shared sparse basis,
while scalar targets specify the score being explained. This separates a
hidden state's target-independent readout profile from support for a particular
token, reference contrast, group contrast, or local margin. The resulting local
account is explicit: each displayed feature term is the product of a target
coefficient fixed by unembedding-row geometry and a context-dependent projection
of the current hidden state, with offset and reconstruction error reported
alongside the feature sum.

Across the model suite, replacement fidelity, target reconstruction, and sign
agreement identify faithful regimes. Within those regimes, SRP shows that
readout explanations are query-local: changing the target changes which active
features count as support, and changing the prompt reweights the same token row
through the hidden-state profile. The case studies expose feature competition in
polysemous token preferences, benchmark-format contrasts, and feature-resolved
DLA, while the readout-edit experiment tests whether selected directions move
held-out lexical scores.

SRP therefore provides a calibrated accounting layer for the final readout. In
faithful regimes, an atomic logit or margin decomposes into sparse output-row
feature support plus measured residual error. These readout-level accounts
connect lens-style summaries to intervention-based mechanistic analysis, while
the reported diagnostics mark where the sparse account is reliable.
